%% kapa: exam=1 round=24 year=2013 period=1 form=A subject=civil_law subject_ko=민법 seq=37 number=37
%% source: Q-Net 감정평가사 기출문제 (cst00309) artlSeq=5073802 file=extracted/1008/5073802/2013년 제24회 감정평가사 1차시험 기출문제/1교시(A형).hwp
%% conversion: kordoc 4.24.51 -> tools/gen1.py (auto); review=proofread-sonnet; flags=none
\begin{question}{37}
분묘기지권에 관한 설명으로 옳지 않은 것은? (다툼이 있으면 판례에 의함)
\begin{choices}
  \item 토지 소유자가 그 토지에 봉분(封墳) 형태로 분묘를 설치한 후 그 봉분의 철거 특약 없이 토지를 타인에게 양도한 경우 분묘기지권이 성립한다.
  \item 부부 중 일방이 먼저 사망하여 이미 분묘가 설치된 후 다른 일방을 쌍분(雙墳) 형태로 합장하여 분묘를 설치하는 것은 그 분묘기지권이 미치는 범위 내라 하더라도 허용되지 않는다.
  \item 분묘기지권의 효력이 미치는 범위 안에서 원래의 분묘를 다른 곳으로 이장하는 것은 허용되지 않는다.
  \item 토지소유자라도 분묘기지권을 침해하는 공작물을 설치할 수 없다.
  \item 타인의 토지 위에 토지 소유자의 승낙없이 분묘를 설치한 후 20년간 평온ㆍ공연하게 분묘기지를 점유한 경우 그 기지에 대한 등기를 함으로써 분묘기지권을 취득한다.
\end{choices}

eviewflag{choices 4-5 overlap in rhwp render; only part verified}
\end{question}
\answer{5}
